\section{Finite energy and lattice texture charge}
\label{sec:lattice-texture-finite-energy}

Consider the cubic spin-$\tfrac12$ source
\[
 H_{\rm spin}=J\sum_{\langle xy\rangle}
 (\mathbf S_x\cdot\mathbf S_y-\tfrac14),\qquad J>0,
\]
and its chosen staggered-field ground $\omega$. The state-selection rule
is an input. Its order vector and nearest-neighbour covariance satisfy
$\omega(\mathbf S_x)=(-1)^x m_{\rm ord}e_1$ and
$C_{ab}=\omega(S_x^aS_y^b)=\operatorname{diag}(a,b,b)$, with
$\kappa=-(a+b)\geq m_{\rm ord}^2/2>0$ and $b\leq0$.
The latter sign also follows from ground positivity: a one-site rotation
by $\pi$ about $e_1$ changes the energy by $-24Jb$.
The spin Hamiltonian is the fixed-spacing strong-coupling limit of the
centered contact source; no physical Lorentz-spin or Dirac-mass
identification is imposed.

For onsite rotations $R_x\in\mathrm{SO}(3)$, put $n_x=R_xe_1$.
The rotated quantum state has the exact excess energy
\begin{align}
 \Delta E
 &=J\sum_{\langle xy\rangle}\operatorname{Tr}
       [C(R_x^{\mathsf T}R_y-I)] \notag\\
 &\geq \frac{J\kappa}{2}
       \sum_{\langle xy\rangle}|n_x-n_y|^2.
 \label{eq:lattice-texture-source-bound}
\end{align}
Indeed, each summand equals
$J\kappa(1-n_x\cdot n_y)+J(-b)(1+R_{11}-R_{22}-R_{33})$,
where $R=R_x^{\mathsf T}R_y$.
For a rotation through $\theta$ about $u$, the last factor is
$2(1-\cos\theta)u_1^2\geq0$.
Each bond excess is at most $J$, because the source bond has spectrum
$\{-J,0\}$ and the reference bond energy is $J(a+2b-1/4)$.
These statements concern the rotation orbit of $\omega$, not energy
minimized over all quantum states with specified one-point functions.

\begin{proposition}[Shell charge and finite energy]
\label{prop:lattice-texture-shell-charge}
Triangulate each integer-radius cubic shell by dividing every unit square
with corners $p_{00},p_{10},p_{01},p_{11}$ into
$(p_{00},p_{10},p_{01})$ and $(p_{11},p_{01},p_{10})$, oriented outwards.
Call a shell admissible when every triangle's three direction vectors
have pairwise nonnegative dot products. Normalized barycentric
interpolation then defines a sphere map of integer degree $q_r$.
For any finite collection of concentric shells,
\begin{equation}
 \Delta E\geq \frac{J\kappa}{2}N_{\rm bad}
       +\pi J\kappa\sum_{r\ {
m admissible}}|q_r|,
 \label{eq:lattice-texture-mixed-shell-bound}
\end{equation}
where $N_{\rm bad}$ counts nonadmissible shells.
Consequently every finite-energy rotation assignment has admissible,
degree-zero shells at all sufficiently large integer radii.
\end{proposition}

\begin{proof}
For an admissible triangle with unit vectors $u,v,w$, its unsigned
spherical area obeys
\[
 A=2\arctan\frac{|\det(u,v,w)|}{1+u\cdot v+v\cdot w+w\cdot u}
   \leq |u-v|^2+|u-w|^2.
\]
The denominator is at least one, and
$|\det(u,v,w)|\leq|v-u||w-u|$.
Both chords on the right are actual source bonds at the triangle's
square-corner apex. The two triangles count each square side once;
the closed surface counts each source edge twice. Thus
$4\pi|q_r|\leq2\sum_{\rm shell}|n_x-n_y|^2$.
A nonadmissible triangle has a chord square greater than two.
If that chord is the square diagonal, its square is at most twice the
sum of the two source-bond squares at the apex. Every bad shell
therefore has source-bond square sum greater than one.
Distinct concentric shells have disjoint source edges, proving
\eqref{eq:lattice-texture-mixed-shell-bound}.
For infinite assignments, onsite rotations define a state on every local
algebra and energy is the extended sum of nonnegative bond excesses.
Finite energy permits only finitely many bad or charged shells.
\end{proof}

This conclusion requires no separate asymptotic smoothness assumption.
For opposite charges whose disjoint shell families each extend through
$R$ radii with degrees $+q$ and $-q$, the cost is at least
$2\pi J\kappa R|q|$. This is a bound on those assigned directions, not
a derived two-particle potential. Zero excess energy with an identity
exterior forces $n_x=e_1$ everywhere by
\eqref{eq:lattice-texture-source-bound} and lattice connectivity.

A continuous rotation path can change a shell degree only through a
triangle whose convex hull contains the origin. At that crossing,
$u\cdot(v+w)\leq0$. Its two source-bond squares therefore sum to at least
four, giving the absolute floor $\Delta E\geq2J\kappa$. This bound gives
neither an activation energy above the initial state nor a lifetime.
For an assignment supported in a cube of side $N$, a controlled erasure
path satisfies
\[
 \max_{\rm path}\Delta E
 \leq\Delta E_{\rm initial}+J(N^2+N+7).
\]
Reset sites lexicographically along continuous quaternion rotations.
A completed prefix cuts at most $N^2+N+1$ bonds; the rotating site touches
at most six more. The bound follows from the bond spectrum. This path
is a specified control, not autonomous Hamiltonian evolution.

The complete two-flavour field-repair instrument gives
$K_{c,r}P=\delta_{r0}P/\sqrt2$ on the bare one-fermion-per-site carrier.
Here $P$ is the carrier projection, $c$ the coin bit and $r$ the raw syndrome.
With every coin and raw record retained, repair-only schedules preserve
any spin state and untouched reference and append independent fair
coins with raw zero. They select no texture. Source hopping can leave
this carrier, so the identity does not describe interleaved full dynamics.
Nor does the energy theorem identify the spin direction with a prescribed
fermion mass texture or exclude defects outside the stated quantum orbit.
Exact degree counts, bond identities and complete branch checks are in
\path{model/lattice_texture_selection/}.
